\section{Solutions for Chapter~\ref{ch:correlators}}

\paragraph{Exercise~\ref{ex:corr-fifteen}.}
\phantomsection
\label{sol:ex:corr-fifteen}
Write the momenta at $(\infty,1,x,0)$ as
$(-a_i,a_j,-a_k,a_\ell)$. Neutrality requires the
multisets $\{i,k\}$ and $\{j,\ell\}$ to agree.
There are three choices with all labels equal. For $i\ne k$,
the choices $(j,\ell)=(i,k)$ and $(k,i)$ give six each.
The finite-distance factors are
$x^{-(a_k,a_\ell)}(1-x)^{-(a_j,a_k)}$.
They are, respectively,
$x^{-4}(1-x)^{-4}$,
$x^{-4}(1-x)^{-2}$, and
$x^{-2}(1-x)^{-4}$.
Each assignment has coefficient $1/9$, giving the three
terms in the chapter.



\paragraph{Exercise~\ref{ex:corr-blocks}.}
\phantomsection
\label{sol:ex:corr-blocks}
At order four the blocks of exchanged weights one, two,
and three contribute $2/3$, $21/5$, and $29/3$.
Subtracting them from the reduced-correlator coefficient gives
\[
 A_4=\frac{65}{3}-\frac23-\frac{21}{5}-\frac{29}{3}
     =\frac{107}{15}.
\]
In the charged ordering the reduced correlator is
$x^4G_+(x)=(4/3)x^6+\cdots$, so the first exchanged
global primary has weight six and squared coefficient $4/3$.
The corresponding unit-normalized three-point coefficient
has magnitude $2/\sqrt3$.



\paragraph{Exercise~\ref{ex:corr-stress}.}
\phantomsection
\label{sol:ex:corr-stress}
Only pure-factor contractions survive, and
$\langle T_i(\infty)T_i(1)T_i(0)\rangle=c_i$.
For $A=c_2T_1-c_1T_2$ the unnormalized three-point
coefficient is $c_1c_2(c_2^2-c_1^2)$.
Divide its square by
$\norm A^6=[c_1c_2(c_1+c_2)/2]^3$ to obtain
\[
 C_{\widehat A\widehat A\widehat A}^2
 =\frac{8(c_2-c_1)^2}{c_1c_2(c_1+c_2)}.
\]
When $c_1=c_2$, exchanging the two Virasoro factors within their
two-stress vacuum subalgebra sends $A$ to $-A$, so its odd
correlator vanishes. Its positive norm
does not vanish; it remains a nonzero field.



\paragraph{Exercise~\ref{ex:corr-spin3}.}
\phantomsection
\label{sol:ex:corr-spin3}
At the Potts point $5c_0+22=26$ and
$\lambda(5c_0+22)=(4/13)26=8$, so the weight-four
primary norm is zero. For the simple pole in the four-point
function, $\frac12\partial b$ contributes
$\frac12(4B)=2B$ and
$-\frac1{12}\partial^3T_0$ contributes
$-72N/12=-6N=-2c_0$.
The total is $2B-2c_0$.



\paragraph{Exercise~\ref{ex:corr-lightgram}.}
\phantomsection
\label{sol:ex:corr-lightgram}
Using $q=(2,0,-2)$ gives
\[
 K=\frac1{39}
 \begin{pmatrix}
 24&-13&-11\\
 -13&26&-13\\
 -11&-13&24
 \end{pmatrix}.
\]
The vectors $(1,1,1)$, $(1,0,-1)$, and $(1,-2,1)$
are eigenvectors with eigenvalues $0$, $35/39$, and $1$.
They are mutually orthogonal and span the space, so $K$
is positive semidefinite. Positivity concerns the quadratic
form, not the signs of its individual matrix entries.


